\section{References}

\begin{frame}{}
    \LARGE \textbf{References}
\end{frame}

\begin{frame}[allowframebreaks]{References \& Further Reading}
    \small

    \begin{thebibliography}{99}

% Key Textbooks
\bibitem{jurafsky2026slp3}
Jurafsky, D., \& Martin, J. H. (2026). \textit{Speech \& Language Processing} (3rd edition draft). \url{https://web.stanford.edu/~jurafsky/slp3/}

\bibitem{rabiner1993fundamentals}
Rabiner, L., \& Juang, B.-H. (1993). \textit{Fundamentals of Speech Recognition}. Prentice Hall.

\bibitem{deng2013paradigms}
Deng, L., \& Li, X. (2013). Machine learning paradigms for speech recognition: An overview. \textit{IEEE Transactions on Audio, Speech, and Language Processing, 21}(5), 1060--1089.

\bibitem{bishop2006pattern}
Bishop, C. M. (2006). \textit{Pattern Recognition and Machine Learning}. Springer.

% Features & Signal Processing
\bibitem{kubichek1993mcd}
Kubichek, R. F. (1993). Mel-cepstral distance measure for objective speech quality assessment. \textit{IEEE Pacific Rim Conference on Communications, Computers and Signal Processing (PACRIM)}, 125--128.

% Seminal Papers (ASR)
\bibitem{rabiner1989tutorial}
Rabiner, L. (1989). A tutorial on hidden Markov models and selected applications in speech recognition. \textit{Proceedings of the IEEE, 77}(2), 257--286.

\bibitem{graves2006ctc}
Graves, A., Fernández, S., Gomez, F., \& Schmidhuber, J. (2006). Connectionist temporal classification: Labelling unsegmented sequence data with recurrent neural networks. \textit{ICML}, 369--376.

\bibitem{hannun2017ctc}
Hannun, A. (2017). Sequence modeling with CTC. \textit{Distill}. \url{https://distill.pub/2017/ctc/}

\bibitem{graves2012rnn}
Graves, A. (2012). Sequence transduction with recurrent neural networks. \textit{arXiv preprint arXiv:1211.3711}.

\bibitem{chan2016las}
Chan, W., Jaitly, N., Le, Q. V., \& Vinyals, O. (2016). Listen, attend and spell: A neural network for large vocabulary conversational speech recognition. \textit{ICASSP}, 4960--4964. arXiv:1508.01211.

\bibitem{park2019specaugment}
Park, D. S., Chan, W., Zhang, Y., Chiu, C.-C., Zoph, B., Cubuk, E. D., \& Le, Q. V. (2019). SpecAugment: A simple data augmentation method for automatic speech recognition. \textit{Interspeech}, 2613--2617. arXiv:1904.08779.

\bibitem{gulati2020conformer}
Gulati, A., Qin, J., Chiu, C.-C., et al. (2020). Conformer: Convolution-augmented transformer for speech recognition. \textit{Interspeech}, 5036--5040. arXiv:2005.08100.

\bibitem{rekesh2023fastconformer}
Rekesh, D., Koluguri, N. R., Kriman, S., et al. (2023). Fast Conformer with linearly scalable attention for efficient speech recognition. \textit{IEEE ASRU}, 1--8. arXiv:2305.05084.

\bibitem{majumdar2021citrinet}
Majumdar, S., Balam, J., Hrinchuk, O., Lavrukhin, V., Noroozi, V., \& Ginsburg, B. (2021). Citrinet: Closing the gap between non-autoregressive and autoregressive end-to-end models for automatic speech recognition. \textit{arXiv preprint arXiv:2104.01721}.

\bibitem{radford2022whisper}
Radford, A., Kim, J. W., Xu, T., Brockman, G., McLeavey, C., \& Sutskever, I. (2022). Robust speech recognition via large-scale weak supervision. \textit{arXiv preprint arXiv:2212.04356}.

% Self-Supervised Speech Representations
\bibitem{baevski2020wav2vec2}
Baevski, A., Zhou, H., Mohamed, A., \& Auli, M. (2020). wav2vec 2.0: A framework for self-supervised learning of speech representations. \textit{NeurIPS 33}. arXiv:2006.11477.

\bibitem{hsu2021hubert}
Hsu, W.-N., Bolte, B., Tsai, Y.-H. H., Lakhotia, K., Salakhutdinov, R., \& Mohamed, A. (2021). HuBERT: Self-supervised speech representation learning by masked prediction of hidden units. \textit{IEEE/ACM Transactions on Audio, Speech, and Language Processing, 29}, 3451--3460. arXiv:2106.07447.

\bibitem{chen2022wavlm}
Chen, S., Wang, C., Chen, Z., et al. (2022). WavLM: Large-scale self-supervised pre-training for full stack speech processing. \textit{IEEE Journal of Selected Topics in Signal Processing, 16}(6), 1505--1518. arXiv:2110.13900.

% Seminal Papers (TTS)
\bibitem{wang2017tacotron}
Wang, Y., et al. (2017). Tacotron: Towards end-to-end speech synthesis. \textit{arXiv preprint arXiv:1703.10135}.

\bibitem{shen2018tacotron2}
Shen, J., Pang, R., Weiss, R. J., et al. (2018). Natural TTS synthesis by conditioning WaveNet on mel spectrogram predictions. \textit{ICASSP}, 4779--4783. arXiv:1712.05884.

\bibitem{oord2016wavenet}
van den Oord, A., et al. (2016). WaveNet: A generative model for raw audio. \textit{arXiv preprint arXiv:1609.03499}.

\bibitem{ren2021fastspeech2}
Ren, Y., Hu, C., Tan, X., Qin, T., Zhao, S., Zhao, Z., \& Liu, T.-Y. (2021). FastSpeech 2: Fast and high-quality end-to-end text to speech. \textit{ICLR}. arXiv:2006.04558.

\bibitem{kong2020hifigan}
Kong, J., Kim, J., \& Bae, J. (2020). HiFi-GAN: Generative adversarial networks for efficient and high fidelity speech synthesis. \textit{NeurIPS 33}. arXiv:2010.05646.

\bibitem{popov2021gradtts}
Popov, V., Vovk, I., Gogoryan, V., Sadekova, T., \& Kudinov, M. (2021). Grad-TTS: A diffusion probabilistic model for text-to-speech. \textit{ICML}, PMLR 139, 8599--8608. arXiv:2105.06337.

\bibitem{kim2021vits}
Kim, J., Kong, J., \& Son, J. (2021). VITS: Conditional variational autoencoder with adversarial learning for end-to-end text-to-speech. \textit{ICML}. arXiv:2106.06103.

% Diarization, Emotion & Visual Speech
\bibitem{park2022diarization}
Park, T. J., Kanda, N., Dimitriadis, D., Han, K. J., Watanabe, S., \& Narayanan, S. (2022). A review of speaker diarization: Recent advances with deep learning. \textit{Computer Speech \& Language, 72}, 101317. arXiv:2101.09624.

\bibitem{lin2020diarization}
Lin, Q., Hou, Y., \& Li, M. (2020). Self-attentive similarity measurement strategies in speaker diarization. \textit{Interspeech}, 284--288.

\bibitem{bredin2020pyannote}
Bredin, H., Yin, R., Coria, J. M., et al. (2020). pyannote.audio: Neural building blocks for speaker diarization. \textit{ICASSP}, 7124--7128.

\bibitem{busso2008iemocap}
Busso, C., Bulut, M., Lee, C.-C., et al. (2008). IEMOCAP: Interactive emotional dyadic motion capture database. \textit{Language Resources and Evaluation, 42}(4), 335--359.

\bibitem{chen2023wav2vec2ser}
Chen, L.-W., \& Rudnicky, A. (2023). Exploring wav2vec 2.0 fine tuning for improved speech emotion recognition. \textit{ICASSP}, 1--5. arXiv:2110.06309.

\bibitem{ma2024emotion2vec}
Ma, Z., Zheng, Z., Ye, J., Li, J., Gao, Z., Zhang, S., \& Chen, X. (2024). emotion2vec: Self-supervised pre-training for speech emotion representation. \textit{Findings of ACL 2024}, 15747--15760. arXiv:2312.15185.

\bibitem{assael2016lipnet}
Assael, Y. M., Shillingford, B., Whiteson, S., \& de Freitas, N. (2016). LipNet: End-to-end sentence-level lipreading. \textit{arXiv preprint arXiv:1611.01599}.

\bibitem{chung2017lrs}
Chung, J. S., Senior, A., Vinyals, O., \& Zisserman, A. (2017). Lip reading sentences in the wild. \textit{CVPR}. arXiv:1611.05358.

% Toolkits
\bibitem{ravanelli2021speechbrain}
Ravanelli, M., et al. (2021). SpeechBrain: A general-purpose speech toolkit. \textit{arXiv preprint arXiv:2106.04624}.

\bibitem{rivadocs}
NVIDIA Riva, ASR models reference. \url{https://docs.nvidia.com/deeplearning/riva/user-guide/docs/reference/models/asr.html}

% Evaluation
\bibitem{itu800}
ITU-T Recommendation P.800 (1996). \textit{Methods for Subjective Determination of Transmission Quality}. International Telecommunication Union, Geneva.

\bibitem{papineni2002bleu}
Papineni, K., Roukos, S., Ward, T., \& Zhu, W.-J. (2002). BLEU: A method for automatic evaluation of machine translation. \textit{ACL}.

% Speech Datasets (ASR)
\bibitem{panayotov2015librispeech}
Panayotov, V., Chen, G., Povey, D., \& Khudanpur, S. (2015). LibriSpeech: An ASR corpus based on public domain audio books. \textit{ICASSP}.

\bibitem{ardila2020commonvoice}
Ardila, R., et al. (2020). Common Voice: A massively-multilingual speech corpus. \textit{LREC}, 4218--4222. arXiv:1912.06670.

\bibitem{hernandez2018tedlium}
Hernandez, F., Nguyen, V., Ghannay, S., Tomashenko, N., \& Estève, Y. (2018). TED-LIUM 3: Twice as much data and corpus repartition for experiments on speaker adaptation. \textit{Speech and Computer (SPECOM)}, LNCS 11096, 198--208. Springer.

\bibitem{carletta2005ami}
Carletta, J., et al. (2005). The AMI meeting corpus: A pre-announcement. \textit{MLMI}, LNCS 3869, 28--39. Springer.

% Speech Datasets (TTS & SLU)
\bibitem{ljspeech}
LJSpeech dataset. \url{https://keithito.com/LJ-Speech-Dataset/}

\bibitem{vctk}
CSTR VCTK Corpus (version 0.92). \url{https://datashare.ed.ac.uk/handle/10283/3443}

\bibitem{hifitts}
Bakhturina, E., Lavrukhin, V., Ginsburg, B., \& Zhang, Y. (2021). Hi-Fi multi-speaker English TTS dataset. \textit{arXiv preprint arXiv:2104.01497}. Download: \url{https://www.openslr.org/109/}

\bibitem{atis}
ATIS intent/slot dataset. \url{https://www.kaggle.com/datasets/siddhadev/ms-cntk-atis}

\bibitem{slurp}
SLURP dataset (audio hosted on Zenodo). \url{https://github.com/pswietojanski/slurp}

% Online Demos
\bibitem{whisperdemo}
Whisper ASR Demo. \url{https://huggingface.co/spaces/openai/whisper}

\bibitem{tacotron2demo}
Tacotron 2 + WaveGlow Demo. \url{https://colab.research.google.com/github/NVIDIA/DeepLearningExamples/blob/master/PyTorch/SpeechSynthesis/Tacotron2/notebooks/Tacotron2.ipynb}

\bibitem{googlets}
Google Cloud Text-to-Speech. \url{https://cloud.google.com/text-to-speech}

\bibitem{coquitts}
Coqui TTS (maintained fork). \url{https://github.com/idiap/coqui-ai-TTS}

% Surveys
\bibitem{li2022asr}
Li, J. (2022). Recent advances in end-to-end automatic speech recognition. \textit{APSIPA Transactions on Signal and Information Processing, 11}(1), e8. arXiv:2111.01690.

\bibitem{prabhavalkar2024e2e}
Prabhavalkar, R., Hori, T., Sainath, T. N., Schlüter, R., \& Watanabe, S. (2024). End-to-end speech recognition: A survey. \textit{IEEE/ACM Transactions on Audio, Speech, and Language Processing, 32}, 325--351. arXiv:2303.03329.

\bibitem{qin2021slu}
Qin, L., Xie, T., Che, W., \& Liu, T. (2021). A survey on spoken language understanding: Recent advances and new frontiers. \textit{IJCAI}, 4577--4584. arXiv:2103.03095.

\bibitem{tan2021tts}
Tan, X., Qin, T., Soong, F., \& Liu, T.-Y. (2021). A survey on neural speech synthesis. \textit{arXiv preprint arXiv:2106.15561}.

\end{thebibliography}

\end{frame}
